\documentclass{article}
\usepackage[utf8]{inputenc}
\usepackage[T1]{fontenc}
\usepackage{geometry}
\geometry{
    a4paper,
    left=2.5cm,
    right=2.5cm,
    top=2.5cm,
    bottom=2.5cm
}
\usepackage{graphicx}
\usepackage{float}
\usepackage{hyperref}
\usepackage[english]{babel}
\usepackage{amsmath}
\usepackage{amssymb}
\usepackage{array}
\usepackage{booktabs}
\usepackage{caption}
\usepackage{subcaption}
\usepackage{enumitem}
\usepackage{fancyhdr}
\usepackage{titlesec}
\usepackage{url}
\usepackage{textcomp}
\usepackage{xcolor}
\usepackage{listings}
\usepackage{multirow}
\usepackage{longtable}
\usepackage{tabularx}
\usepackage{times}

% Set up page style
\pagestyle{fancy}
\fancyhf{}
\rfoot{\thepage}
\renewcommand{\headrulewidth}{0pt}

% Title formatting
\titleformat{\chapter}[display]{\normalfont\bfseries}{\chaptertitle\thechapter}{0pt}{\Huge}
\titlespacing*{\chapter}{0pt}{50pt}{40pt}

% Document info
\title{Template Document}
\author{Academic Support Agent}
\date{\today}

\begin{document}
\maketitle

\maketitle


\begin{frame}
    \frametitle{Conte do}
    
\end{frame}

\section{Introdu o}


\begin{frame}
    \frametitle{Sobre \emph{Beamer}}
\begin{itemize}
    \item Em $LaTex$, o $Beamer$ serve para fazer apresenta es;
    \item  Diferente de programas $WYSWYG$;
    \item  Uma apresenta o $Beamer$   como qualquer outro documento LaTeX, cont m:
    \begin{itemize}
        \item Pre mbulo e um corpo;
        \item O pre mbulo pode-se dizer que   o `` ndice'', tipo do documento e pacotes;
        \item O corpo cont m $sections$ e $subsections$;
        \item Os dispositivos dever o ser estruturados utilizando ambientes de $item$ e $enumerate$, ou texto simples (curto).
\end{itemize} 
\end{itemize}
\end{frame}

\section{Estrutura do Documento}
\subsection{Exemplos}

\begin{frame}
\frametitle{Simples Exemplo}
\scriptsize{
\textbackslash documentclass\{beamer\}\\ \vspace{0.1cm}

\textbackslash title\{\textbackslash Apresenta o IPB atrav s em \LaTeX\{\} \emph{Beamer}\} \\
\textbackslash author\{Primeiro Autor\}\\
\textbackslash institute\{Licentura em Curso \textbackslash\textbackslash Instituto Polit cnico de Bragan a\}\\
\textbackslash date\{\textbackslash scriptsize\{\textbackslash today\}\}
 \\ \vspace{0.3cm}

\textbackslash begin\{document\} \\ \vspace{0.1cm}

\hspace{0.5cm}\textbackslash maketitle\\ \vspace{0.1cm}

\hspace{0.5cm}\textbackslash section\{Nome da $Section$\}\\
\hspace{0.5cm}\textbackslash begin\{frame\}\\
\hspace{1cm}\textbackslash frametitle\{Nome da $Frame$\}\\
\hspace{1cm}\textbackslash Here is one slide.\\
\hspace{0.5cm}\textbackslash end\{frame\}\\ \vspace{0.1cm}

\textbackslash end\{document\}\\
}

\end{frame}


\subsection{Corpo}

\begin{frame}
\frametitle{T tulo da \emph{frame}}
Usar os seguintes comandos no pre mbulo:
\begin{block}{Comando para o t tulo da $frame$}
\scriptsize{\tt{
\textbackslash title\{T tulo\} \\
\textbackslash subtitle\{Sub T tulo\} \\
\textbackslash author\{Nome\}\\
\textbackslash institute\{Nome do Curso\}\\
\textbackslash date\{Data\} \\
}}
\end{block}
e o comando \textbackslash \tt{maketitle}, 
\begin{block}{Gerar a $frame$ no corpo do documento}
\scriptsize{\tt{
\textbackslash begin\{document\} \\ 
\textbackslash maketitle\\ 
\vdots
\textbackslash end\{document\}\\
}}
\end{block}
\end{frame}

\begin{frame}
\frametitle{Criar \emph{frame}}
Para criar $frames$ utiliza-se os seguintes comandos:
\begin{block}{}
\scriptsize{\tt{
\textbackslash begin\{frame\} \\
\textbackslash frametitle\{...\} \\
\textbackslash framesubtitle\{...\} \\
...\\
\textbackslash end\{frame\}
}}
\end{block}
Ou, 
\begin{block}{}
\scriptsize{\tt{
\textbackslash frame\{ \\
\textbackslash frametitle\{...\}\\
\textbackslash framesubtitle\{...\}\\
...\\
\}
}}
\end{block}
\end{frame}

\begin{frame}
\frametitle{\emph{section} e \emph{subsections}}
\vspace{1cm}
\scriptsize{\tt{
...\\
\textbackslash end\{frame\}\\
\vspace{0.5cm}
\alert{
\textbackslash section\{section name\}\\
\textbackslash subsection\{subsection name\}\\
\textbackslash subsubsection\{subsubsection name\}}\\
\vspace{0.5cm}
\textbackslash begin\{frame\}\\
...\\
}
}
\end{frame}

\begin{frame}
\frametitle{Conte do} 

\begin{block}{Manual}
\scriptsize{\tt{
...\\
\textbackslash end\{frame\}\\
\vspace{0.1cm}
\textbackslash begin\{frame\}\\
\textbackslash frametitle\{Lista de Conte do\}\\
\textbackslash tableofcontents[currentsection]\\
\textbackslash end\{frame\}\\
\vspace{0.2cm}
\textbackslash begin\{frame\}\\
...\\
}}
\end{block}

\begin{block}{Autom tico}
\scriptsize{\tt{
...\\
\textbackslash AtBeginSection[]\{\\
\textbackslash begin\{frame\}\{Outline\}\\
\textbackslash tableofcontents[currentsection]\\
\textbackslash end\{frame\}\}\\
...\\
}}
\end{block}
\end{frame}

\section{Estratura da \emph{frame}}
\subsection{Criar colunas}


%%%%%%%%%%%%%%%%%%%%%%%%%%%%%%Frame 15
\begin{frame}
\frametitle{Duas $colums$}
\begin{block}{Exemplo de duas $colums$}\scriptsize{
\textbackslash begin\{columns\}\\
\textbackslash column\{.4\textbackslash textwidth\}\\
Left column\\
\textbackslash column\{.4\textbackslash textwidth\}\\
Right column\\
\textbackslash end\{columns\}\\
}
\end{block}
\vspace{1cm}
%\onslide<2>{
\begin{columns}
\column{.4\textwidth}
Left column\\
\column{.4\textwidth}
Right column\\
\end{columns}
%}
\end{frame}

\subsection{Criar \emph{Block's}}
\begin{frame}
\frametitle{Block}

\begin{block}{Beamer Introduction}
Beamer is a { \LaTeX} class.
\end{block}

\end{frame}

\subsection{Criar listas}
\begin{frame}[fragile]
\frametitle{\emph{itemize}}

\begin{columns}
\column{.5\textwidth}
\begin{block}
\scriptsize{
\begin{verbatim}
\begin{itemize}
\item The first one.
\item The second one.
\begin{itemize}
\item The larger one.
\item The smaller one.
\end{itemize}
\item The third one.
\end{itemize}
\end{verbatim}
}
\end{block}
\column{.5\textwidth}

\begin{itemize}
\item The first one.
\item The second one.
\begin{itemize}
\item The larger one.
\item The smaller one.
\end{itemize}
\item The third one.
\end{itemize}
\end{columns}
\vspace{1.5cm}
\href{https://tug.ctan.org/macros/latex/contrib/beamer/doc/beameruserguide.pdf}{Clique aqui para mais informa es.}
\end{frame}

\begin{frame}[fragile]
\frametitle{\emph{enumerate}}

\begin{columns}
\column{.5\textwidth}
\begin{block}
\scriptsize{
\begin{verbatim}
\begin{enumerate}
\item The first one.
\item The second one.
\begin{enumerate}
\item The large one.
\item The small one.
\end{enumerate}
\item The third one.
\end{enumerate}
\end{verbatim}
}
\end{block}
\column{.5\textwidth}

\begin{enumerate}
\item The first one.
\item The second one.
\begin{enumerate}
\item The large one.
\item The small one.
\end{enumerate}
\item The third one.
\end{enumerate}
\end{columns}
\vspace{1.5cm}
\href{https://tug.ctan.org/macros/latex/contrib/beamer/doc/beameruserguide.pdf}{Clique aqui para mais informa es.}

\end{frame}
%\section{Refer ncias}
%\begin{frame}[allowframebreaks]
 %   
    %
  %  \nocite{*} % used here because no citation happens in slides
    % if there are too many try use 
   %  \tiny
%\end{frame}
\end{document}